% ICLR 2027: these statements are outside the main-text page limit.
% Headings follow the official template; factual disclosures follow the study record.
\subsection*{AI use statement}
The discovery campaigns used \texttt{gemini-3.5-flash} and
\texttt{gemini-3.1-pro-preview} through AlphaEvolve; the controlled study
used Gemini 3.5 Flash with a local controller. These models evolved
search programs using researcher proposals and exact evaluation feedback
(Sections~\ref{sec:pipeline} and~\ref{sec:ablation}). LLM-based coding
assistants, including Codex, supported research design, implementation,
mathematical reasoning---including the coordinate-normalization argument
in Lemma~\ref{lem:coordinate-normalization}---and result analysis.
AI also assisted literature review, language editing, and preparation
of the manuscript and supporting artifacts. Mathematical and computational
claims were checked using the verification procedures in
Section~\ref{sec:refute} and Appendix~\ref{app:lean}.
The authors take responsibility for the final text, mathematical claims,
proofs, code and reported results, including all AI-assisted contributions.

\subsection*{Ethics statement}
This work studies algebraic code constructions and search programs;
it does not involve experiments on human participants or use
participant-level or personal datasets. Its potential benefit is to
improve methods for quantum error correction and communication, while
the reported idealized parameters do not establish hardware-level
reliability or operational security. Research-integrity considerations
are addressed through dated sources, explicit provenance, independent
verification, and separation of new parameter claims from retrospective
reconstruction. Unsuccessful searches are distinguished from proved
nonexistence. The controlled-study records preserve failed and excluded
attempts, unavailable responses and protocol amendments, and its reviewer
bundle omits account credentials and machine-local configuration.
Researcher and AI contributions are disclosed, external sources are
cited, and the artifact retains software licensing information. Compute
allowances and model usage are documented in Appendices~\ref{app:campaign}
and~\ref{app:ablations}.

\subsection*{Reproducibility statement}
The artifact supports reconstruction and independent checking of the
mathematical outputs: generator matrices specify the finite codes,
implementations generate family instances, and Appendix~\ref{app:theorem}
provides the general construction and proof of Theorem~\ref{thm:unified}.
The initial and later task specifications in Appendix~\ref{app:prompts}
document the transition from initial search directions to reuse of
the researcher-derived qubit family. Preserved programs can be replayed with fixed
seeds and evaluation budgets as described in Section~\ref{sec:repro}
and the \artifactlink{}. Appendix~\ref{app:ablations} and the accompanying
study records supply generation requests, selections and held-out
outcomes for the separate controlled comparison. Computational checks
verify finite cases; general validity rests on the proof.
Appendix~\ref{app:lean} gives the exact formal goal and proof structure
for Theorem~\ref{thm:unified}. The artifact includes the Lean sources,
pinned dependencies and separate kernel and Comparator verification records.
Fresh LLM sampling is not part of this replay. External solver
and Magma checks are reported separately from completed Python checks.
The artifact is released under the MIT license and is available at
\url{https://anonymous.4open.science/r/llm_guided_EAQECC_search-4099}.
